% !PRINTED 167
\weberpage{167}

%<<CONT
Theilers auftreten, nur dass hier die unzerlegbaren Functionen
die Rolle der Primzahlen übernehmen.

Wir führen noch einen solchen Satz an:

\gesperrt{IV. Sind $u$, $v$ relativ prim und $Uu$, $Uv$ durch $w$
theilbar, so ist $U$ durch $w$ theilbar.}

Denn zerlegt man die ganzen rationalen Functionen $U, u, v, w$
in ihre unzerlegbaren Factoren, so muss irgend ein Factor von
$w$, da er nicht in $u$ und $v$ zugleich vorkommen kann, in $U$ auf\-
gehen. Hebt man ihn aus $U$ und $w$ weg, so kann man ebenso
für einen nächsten Factor von $w$ schliessen u. s. f.

Wir betrachten nun ganze rationale Functionen einer Ver\-
änderlichen $t$
\[
f(t) = u_0 t^m + u_1 t^{m-1} + \cdots + u_{m-1} t + u_m,
\]
deren Co\"efficienten ganze rationale Functionen von $n$ Veränder\-
lichen $x$ sind, von denen $t$ unabhängig ist.

Sind die Co\"efficienten $u_0, u_1 \ldots u_m$ ohne gemeinsamen
% [corrected 2026-08-07, 8x zoom strip 4/8] Two terms letterspaced at their definitions:
% "primitiv", and "der Theiler der Function". The run STOPS before $f(t)$ -- "f(t) genannt
% (vgl. §. 2)." is set normally -- so it is one word narrower than the query claimed. Claimed by
% the translation pass and independently by the sperrung pass on a different model.
Theiler, so heisst $f(t)$ \gesperrt{primitiv}, im anderen Falle wird der
grösste gemeinschaftliche Theiler der Co\"efficienten $u_0, u_1 \ldots u_m$
\gesperrt{der Theiler der Function} $f(t)$ genannt (vgl. §.\,2).

Wir schliessen, immer unter Voraussetzung der Gültigkeit
von I:

\gesperrt{V. Das Product von zwei primitiven Functionen $f(t)$
und $\varphi(t)$ ist wieder eine primitive Function und
der Theiler eines Productes zweier imprimitiver
Functionen ist gleich dem Product der Theiler
beider Factoren.}

Der Beweis ist ganz so wie für den entsprechenden Satz
in §.\,2.

Es seien
\weq{1}{
\begin{aligned}
f(t) &= u_0 t^m + u_1 t^{m-1} + \cdots + u_m \\
\varphi(t) &= v_0 t^\mu + v_1 t^{\mu-1} + \cdots + v_\mu
\end{aligned}
}
zwei primitive Functionen und
\weq{2}{
F(t) = U_0 t^{m+\mu} + U_1 t^{m+\mu-1} + \cdots + U_{m+\mu}
}
ihr Product. Es ist dann, wenn $r$ und $s$ irgend zwei Ziffern aus
den Reihen $0, 1 \ldots m$, und $0, 1, 2 \ldots \mu$ bedeuten (§.\,2)
\weq{3}{
\begin{aligned}
U_{r+s} &= u_r v_s + u_{r-1} v_{s+1} + \cdots \\
&\qquad + u_{r+1} v_{s-1} + \cdots
\end{aligned}
}
%CONT>>
